\subsection{مبانی نظری و فرمول‌بندی ریاضی}
فرض کنید تانسورهای پرس‌وجو، کلید و مقدار پس از تفکیک سرها و اعمال احتمالی تعبیه‌های چرخشی مکانی (\lr{RoPE}) به صورت $\mathbf{Q} \in \mathbb{R}^{B \times H \times L \times d}$ و $\mathbf{K}, \mathbf{V} \in \mathbb{R}^{B \times H \times N_k \times d}$ باشند. 

در پردازش دوطرفه انکودر (\lr{Bidirectional})، میانگین فضایی کلیدها با در نظر گرفتن ماسک توکن‌های معتبر محاسبه و کسر می‌گردد:
\begin{equation}
\boldsymbol{\mu}_k = \frac{1}{\sum_{j=1}^{N_k} m_j} \sum_{j=1}^{N_k} \mathbf{K}_j \cdot m_j, \quad \tilde{\mathbf{K}} = \mathbf{K} - \boldsymbol{\mu}_k
\end{equation}
سپس نگاشت غیرخطی نامنفی جهت تضمین مثبت‌بودن شباهت به کار می‌رود:
\begin{equation}
\hat{\mathbf{Q}} = \operatorname{ReLU}(\mathbf{Q}) + \epsilon, \quad \hat{\mathbf{K}} = \operatorname{ReLU}(\tilde{\mathbf{K}}) + \epsilon
\end{equation}
به واسطه خاصیت خطی، ماتریس تجمعی کلید-مقدار با پیچیدگی $\mathcal{O}(N \cdot d^2)$ تشکیل می‌شود:
\begin{equation}
\mathbf{S}_{\text{kv}} = \hat{\mathbf{K}}^\top \mathbf{V} \in \mathbb{R}^{B \times H \times d \times d}, \quad \mathbf{z} = \sum_{j=1}^{N_k} \hat{\mathbf{K}}_j^\top
\end{equation}
و خروجی توجه با ضرب برداری و نرمال‌سازی به دست می‌آید:
\begin{equation}
\mathbf{O} = \frac{\hat{\mathbf{Q}} \mathbf{S}_{\text{kv}}}{\hat{\mathbf{Q}} \mathbf{z} + \epsilon}
\end{equation}

در پردازش علّی دیکودر (\lr{Causal Decoding})، میانگین متحرک توکن‌های گذشته در هر گام زمانی به صورت بازگشتی یا با استفاده از جمع تجمعی محاسبه می‌شود تا از نشت اطلاعات آینده جلوگیری به عمل آید:
\begin{equation}
\boldsymbol{\mu}_{k, t} = \frac{1}{t} \sum_{\tau=1}^t \mathbf{K}_\tau, \quad \tilde{\mathbf{K}}_t = \mathbf{K}_t - \boldsymbol{\mu}_{k, t}
\end{equation}
و حالت بازگشتی پنهان با رابطه زیر به‌روزرسانی می‌گردد:
\begin{equation}
\mathbf{S}_t = \mathbf{S}_{t-1} + \hat{\mathbf{K}}_t^\top \mathbf{V}_t, \quad \mathbf{O}_t = \frac{\hat{\mathbf{Q}}_t \mathbf{S}_t}{\hat{\mathbf{Q}}_t \sum_{\tau=1}^t \hat{\mathbf{K}}_\tau^\top + \epsilon}
\end{equation}

که در آن $\sigma_i(\mathbf{W})$ بیانگر $i$-امین مقدار منفرد مرتب‌شده ماتریس به صورت نزولی است. نسبت حفظ انرژی طیفی $\eta_r = \frac{\sum_{i=1}^r \sigma_i^2}{\sum_{i=1}^{d} \sigma_i^2}$ تضمین می‌کند که با انتخاب آستانه $\tau \ge 0.90$، بیش از نود درصد از پراکندگی واریانس داده‌ها حفظ شده و افت سیگنال در انتشار گرادیان رو به عقب به حداقل ممکن کاهش می‌یابد.

\paragraph{تحلیل مجانبی و پیچیدگی اشتراک حافظه سخت‌افزاری}
از منظر پیچیدگی محاسباتی و انتقال داده‌ها در سلسومراتب حافظه (حافظه سراسری \lr{HBM} به حافظه محلی \lr{SRAM})، عملیات ضرب ماتریسی در این مکانیزم نیازمند تخصیص بافر موقت با ابعاد فاکتورگیری‌شده است. برای پردازش دنباله‌ای به طول $N$ با بعد مخفی $D$، بار محاسباتی هسته از فرم درجه دو به رابطه زیر تقلیل می‌یابد:
\begin{equation}
\label{eq:complexity_zeros_attention}
\mathcal{C}_{\text{zeros\_attention}} = \mathcal{O}\left( B \cdot N \cdot D \cdot r + B \cdot H \cdot N \cdot d_k \right) \ll \mathcal{O}\left( B \cdot H \cdot N^2 \cdot d_k \right)
\end{equation}
این بهبود مرتبه پیچیدگی، با کاهش چشمگیر حجم تبادل تانسورها در گذرگاه‌های ارتباطی کارت گرافیک، مانع از اشباع پهنای باند حافظه (پدیده \lr{Memory-Bound}) در فرآیند استنتاج طولانی‌مدت می‌گردد.
